% ---------------------------------------------------------------- Table S2
% Viewers compared on ten capabilities.  Every mark comes from the tool's papers, documentation or code
% (checked 2026-10-05); the notes give the reason for every partial mark and for the less obvious others.
\clearpage
\setlength{\LTleft}{0pt}\setlength{\LTright}{0pt plus 1fill}
\footnotesize
\newcommand{\N}{\mNative{}}\newcommand{\Pt}{\mPartial{}}\newcommand{\X}{\mNone{}}
\begin{longtable}{@{}l*{10}{c}@{}}
\caption[Interactive viewers for 3-D genome data]{\textbf{Interactive viewers for 3-D genome data.}
Fourteen published viewers and U-Chrom rated on ten capabilities, from each tool's papers, documentation and code (checked 2026-10-05).
\mNative{}~native: built in and documented; \mPartial{}~partial: limited as stated in the notes below the table; \mNone{}~not found in the papers, documentation or code.
Columns and rules follow the table.}
\label{tab:S2}\\
\toprule
Viewer & \rot{bulk contact maps} & \rot{per-cell contact maps (scHi-C)} & \rot{per-spot maps (spatial Hi-C)} & \rot{chromatin tracing (many traces)} & \rot{3-D structures} & \rot{genomic tracks} & \rot{cell embedding or table} & \rot{cell or spot positions} & \rot{views linked by locus} & \rot{views linked by cell selection} \\
\midrule
\endfirsthead
\toprule
Viewer & \multicolumn{10}{c}{(continued)}\\
\midrule
\endhead
\multicolumn{11}{@{}l}{\emph{Contact-map browsers}}\\
HiGlass~\citep{higlass} & \N & \X & \X & \X & \X & \N & \X & \X & \N & \X \\
Juicebox.js~\citep{juiceboxjs} & \N & \X & \X & \X & \X & \N & \X & \X & \N & \X \\
3D Genome Browser 2.0~\citep{threedgb2} & \N & \Pt & \X & \X & \X & \N & \X & \X & \N & \X \\
WashU Epigenome Browser~\citep{washu2022,washu3d,washu2025} & \N & \Pt & \X & \X & \N & \N & \X & \X & \N & \X \\[2pt]
\multicolumn{11}{@{}l}{\emph{3-D structure viewers}}\\
HiCube~\citep{hicube} & \N & \X & \X & \X & \N & \N & \X & \X & \N & \X \\
4D Genome Browser~\citep{fourdgb} & \N & \X & \X & \X & \N & \Pt & \X & \X & \N & \X \\
ChromPolymerDB~\citep{chrompolymerdb} & \N & \X & \X & \X & \N & \N & \X & \X & \Pt & \X \\
ChromoSkein (preprint)~\citep{chromoskein} & \X & \X & \X & \X & \N & \N & \X & \X & \N & \X \\[2pt]
\multicolumn{11}{@{}l}{\emph{Chromatin-tracing viewers}}\\
Spacewalk~\citep{spacewalk} & \N & \X & \X & \N & \N & \N & \X & \X & \Pt & \X \\
OLIVE~\citep{olive} & \X & \X & \X & \N & \N & \X & \X & \X & \Pt & \X \\[2pt]
\multicolumn{11}{@{}l}{\emph{Multimodal and single-cell viewers}}\\
Nucleome Browser~\citep{nucleomebrowser} & \N & \Pt & \X & \Pt & \N & \N & \X & \Pt & \N & \X \\
Vitessce~\citep{vitessce} & \Pt & \X & \X & \X & \X & \N & \N & \N & \Pt & \Pt \\
SGS genome browser~\citep{sgs} & \N & \Pt & \X & \X & \X & \N & \N & \N & \N & \Pt \\
Higashi visualization~\citep{higashi} & \Pt & \N & \X & \X & \X & \X & \N & \X & \Pt & \N \\[2pt]
\midrule
U-Chrom (this work) & \N & \N & \N & \N & \N & \N & \N & \N & \Pt & \N \\
\bottomrule
\end{longtable}

\normalsize
\noindent\textbf{Columns.}
\emph{Bulk contact maps}: Hi-C-type contact matrices from files such as .hic, .cool or .mcool.
\emph{Per-cell contact maps}: the maps of many single cells held together, with a cell or a group of cells chosen to show its map or the sum of the group's maps.
\emph{Per-spot maps}: the contact maps of the spots of a tissue section, tied to the spots' positions.
\emph{Chromatin tracing}: imaged 3-D traces of many cells loaded as data.
\emph{3-D structures}: chromatin structures or traces drawn in 3-D.
\emph{Genomic tracks}: 1-D signals along the genome.
\emph{Cell embedding or table}: a cell embedding (such as a UMAP) or a cell table in which cells can be selected.
\emph{Cell or spot positions}: cells or spots drawn at their positions in a tissue section, field of view or image.
\emph{Views linked by locus}: navigating the genome in one view moves the other views (maps, tracks and 3-D).
\emph{Views linked by cell selection}: selecting cells in an embedding or in space re-scopes the contact maps, tracks or 3-D views to those cells.

\medskip\noindent\textbf{Rules.}
(1)~Each tool is rated on its own; a feature of a separate application built around it counts for that application (Vitessce embeds HiGlass, Spacewalk embeds juicebox.js and igv.js), and a viewer built only for one figure of a paper counts as partial.
(2)~Per-cell contact maps count only when a paper or the documentation shows per-cell or per-group maps of a single-cell data set; a statement that a file format could hold them does not.
(3)~A map computed by the viewer from 3-D coordinates is a distance or contact map of those coordinates, not a contact map in the first column.
(4)~Linking by cell selection needs a cell embedding, cell table or spatial view to select in; stepping through traces, alleles or models by number is noted but not counted.
(5)~Where two readings were possible, the stricter mark was given and the note says so.

\medskip\noindent\textbf{Notes.}
\emph{HiGlass}: higlass.js 2.3.5; no cell, trace or 3-D data type in the core or the official plugins.
\emph{Juicebox.js}: v4.9.0; its paper notes that a .hic file can hold a single cell's map, but the viewer has no cell collection or cell selection (rule~2).
\emph{3D Genome Browser 2.0}: the single-cell data are cell-type pseudo-bulk maps computed in advance and loaded as ordinary .mcool data sets; per-cell maps are named as future work; no 3-D rendering despite the name.
\emph{WashU Epigenome Browser}: a dynamic Hi-C track steps through a fixed set of maps, also of single cells, without selecting or summing cells (partial); 3-D models (g3d, also of several cells) are drawn with 3Dmol.js and stay in step with the browser region; the image track shows microscopy images by genomic locus, not cell positions.
\emph{HiCube}: two cases side by side; the 3-D view shows one cell's model of a g3d file.
\emph{4D Genome Browser}: two structures from polymer simulation of bulk Hi-C; tracks are per-bead values painted on the structure and plotted for a selection only (partial).
\emph{ChromPolymerDB}: its single-cell structures are inferred from bulk Hi-C, not measured; one locus query sets all views (partial).
\emph{ChromoSkein}: arXiv preprint only; its 2-D map is a distance map of the 3-D model (rule~3); tracks from BED files.
\emph{Spacewalk}: release r15.33, no paper; contact and distance maps of the loaded traces are computed live; one trace is drawn at a time and picked by number (rule~4); hovering a track or the map highlights the region in 3-D, but all panels stay on the locus of the loaded traces (partial).
\emph{OLIVE}: one allele at a time, chosen from lists by field of view and allele; clicking the distance map highlights a readout pair in 3-D (partial).
\emph{Nucleome Browser}: eight single-cell Hi-C maps were shown as separate tracks (partial); chromatin traces and microscopy images of a chosen cell were shown in viewers built for the paper (partial); views share the viewed and highlighted genomic regions; the core platform passes no cell selection between views.
\emph{Vitessce}: v4.0.10; contact maps only through an undocumented generic HiGlass view (partial); genomic profiles are computed in advance per cell set, and a selection colours them without re-computing them (partial); only the HiGlass-based views follow the genomic position (partial); 3-D views are for image volumes and segmentations.
\emph{SGS genome browser}: Hi-C as .hic tracks computed per cell type (partial); a cell selection re-scopes the single-cell ATAC tracks but not the contact maps (partial).
\emph{Higashi visualization}: reads Higashi's own output only; all cells give a pseudo-bulk but no bulk file can be loaded (partial); a selection in the embedding shows the summed maps of the selected cells; map panels share their ranges, set by bin sliders (partial).
\emph{U-Chrom}: tracks are the signals of the store (per spot, per locus, cell~$\times$~bin matrices), not bigWig files; the Genome, Matrix and Stats views share the locus, but the 3-D view does not follow it (a region is highlighted there separately), hence partial.

\medskip\noindent\textbf{Checked and not listed.}
Tools with no interactive view (Fast-Higashi, scHiCTools, pyHiM), libraries and grammars without an application (Uchimata, the Gosling 3-D grammar), tools without cell-level data that duplicate rows above (4DN Data Portal, which uses HiGlass; CSynth, GenomeFlow, Delta, HiBrowser), and unpublished code (napari-nucleome).
The spatial Hi-C studies whose data are in the atlas (Chen et al.\ 2026, Wang et al.\ 2026, Guo et al.\ 2026) do not provide an interactive viewer.
